\subsection{西州与辽东：早唐的边线从不共享一套成本}
640年唐灭高昌、置西州，\eventref{ST-640-GAOCHANG}{高昌归唐}。长安的诏令将绿洲编进州县名称，吐鲁番文书却提醒我们，田宅、户籍和语言必须由当地人重新写入新的行政关系。军队跨过河西走廊后，驻军要吃粮，官员要找得到翻译、书吏和担保人；一个州的成立并不是当地社会同时变成关中的复制品。出土文书只能照亮有限地点，不能替整个西域作结论，却足以让“直接经营”显出它依赖的地方中介。

645年唐军在辽东久攻安市城不克而撤回，\eventref{ST-645-LIAODONG-WITHDRAWAL}{辽东撤军}。中韩材料都保存了战事的轮廓，守城英雄和伤亡数字则各有叙事层累。西州的经营并没有让东北远征变得容易：辽河、城堡、寒冷和补给迫使军队回头。早唐同时拥有向西设置州县、向北安置突厥部众、向东出征半岛的能力，却不能以同样的粮道和地方关系支撑每一方向。把这种不同写出来，才不会把“贞观边功”误读成一张没有摩擦的扩张图。

663年白江口海战和668年高句丽覆亡使唐、新罗、倭与半岛居民被卷入新的接收。\eventref{ST-663-BAEKGANG}{白江口}的舰船数量、伤亡与参战目的在中日韩文献中并不相同；\eventref{ST-668-GOGURYEO-CONQUEST}{安东都护府}的名号也不证明唐对所有城镇有同等控制。海峡、港口与城堡要求各方继续供给、交涉和选择。唐的国家线在此并未停止，却必须承认它的边疆由别人的道路和政权共同塑成。
